Several limitations of the preserved evidence need to be stated explicitly. The most important is that all final checkpoints were selected by validation accuracy. Held-out test metrics were recomputed for the final CIFAR-10, CIFAR-100, and ModelNet40 checkpoints, but those test metrics were not used for model selection and are still single-run measurements. They should therefore not be compared casually to published Swin results reported under standardized multi-run or carefully tuned benchmark protocols.

The second major limitation is the absence of repeated-seed experiments, error bars, and statistical significance tests. The project supports reproducibility in the engineering sense through configuration files, saved metrics, queue logs, Hugging Face weight exports, and deterministic artifact paths. However, reproducibility in the stronger statistical sense would require rerunning the same configuration multiple times and quantifying variation across seeds or training stochasticity. Those measurements are not available in the preserved artifact set.

The third limitation concerns the gap between sweep-best and final-retrain performance. CIFAR-10 and CIFAR-100 improved substantially after full retraining, while ModelNet40 restored a weaker checkpoint than the best short-budget sweep trial. The final CIFAR-10 and ModelNet40 retraining configurations also lowered drop-path regularization relative to their sweep-selected artifacts, so the report treats the runs as comparable evidence for dimensional generality, not as a perfectly controlled benchmark suite. Better retraining policies, stronger schedule transfer from sweep to full run, or repeated validation of the selected artifact would strengthen the empirical story.

Another limitation is the scope of the empirical comparison itself. The report evaluates two small 2D image datasets and one voxelized 3D object-recognition dataset. This is enough to demonstrate dimensional generality, but not to claim broad superiority across datasets, tasks, or modalities. The ModelNet40 experiment also depends on a voxelized export pipeline rather than on the raw representation used in some other 3D learning pipelines. Its validation split is train-derived and documented by the preserved dataset manifest, so the present report treats the 3D result as evidence that the implementation works on volumetric data, not as a universal statement about 3D vision performance.

With those boundaries in place, the overall conclusion is positive. The project successfully implemented a reusable N-dimensional shifted-window transformer framework in JAX and demonstrated that it can be trained in a practical workflow on CIFAR-10, CIFAR-100, and ModelNet40 voxels. The architecture-level generalization is real in code, not merely conceptual: patch embedding, window partitioning, shifting, relative position handling, and stage composition all operate on configurable spatial dimensionality. The surrounding experimentation framework is likewise functional, supporting sweeps, queue execution, checkpointing, structured metric export, and publishable model-weight bundles.

In that sense, the work fulfills the core objective of a practical AI project based on the original Swin paper. It translates a strong published idea into a broader implementation setting, exposes the engineering decisions needed to make that translation work, and produces valid empirical evidence that the resulting system can learn on multiple data modalities. Future work should prioritize repeated-seed uncertainty estimates, stronger retraining policies after sweep selection, direct comparison to published Swin baselines, and larger-scale datasets where transformer architectures are known to be more competitive. Those additions would strengthen the scientific completeness of the evidence, but they are not required to support the central conclusion of this report: NDSwin-JAX is a credible, traceable N-dimensional implementation of shifted-window transformers in JAX.
